\section{Methodology}
\label{sec:methodology}

\subsection{Problem Definition and Notation}

We formulate drug-target interaction prediction as a binary classification problem. Given a drug $D$, a target protein $T$, and an interaction label $y\in\{0,1\}$, the model estimates $p(y=1\mid D,T)$. Each drug is represented by SMILES-level, graph-level, and conformer-level features, while each target is represented by protein sequence tokens. The goal is to learn drug and target representations that preserve modality-specific structure and support explicit pairwise matching.

\subsection{Overall Framework}

TAMR-DTI uses a multimodal drug-target architecture with three drug views and one protein sequence stream. A drug encoder extracts ChemBERTa SMILES features, GCN topology features, and EGNN 3D conformer features. A target encoder extracts ProtBERT features followed by a local-context protein encoder. The proposed components are inserted at two controlled points: target-aware conformer modulation in the drug/protein representation stage, and a protein-only Mamba residual enhancer before BiIntention. The final interaction representation is produced by BiIntention and passed to an MLP prediction head.

\subsection{Multimodal Drug Representation Learning}

The drug branch uses three complementary views. ChemBERTa provides SMILES-level semantic features, GCN encodes molecular topology, and EGNN encodes coordinate-sensitive conformer features. These features are projected into a shared hidden space before interaction modeling.

\subsection{Target-Aware Conformer Modulation}

The first modification makes conformer aggregation conditional on the target protein. Instead of pooling drug conformers without target context, the model uses protein information to guide conformer selection and produces a target-aware 3D drug representation. This design matches the DTI setting: the same molecule may expose different relevant geometric cues under different protein contexts.

\subsection{Target Representation Learning}

The protein branch starts from ProtBERT token features and applies a local-context protein encoder. We further condition protein token encoding on the fused ligand representation through a FiLM-style modulation path. This ligand-conditioned protein representation is intended to make target encoding sensitive to the paired drug while still preserving the protein sequence stream for later interaction modeling.

\subsection{Mamba-based Protein Representation Refinement}

The second modification inserts a Mamba residual refinement block only on the protein token stream before BiIntention:
\[
  \tilde{P} = P + g \cdot \mathrm{Mamba}(P),
\]
where $P$ denotes protein tokens and $g$ is a learned residual gate initialized from a negative bias. The best current setting uses a gate bias of $-2$, corresponding to an initial injection strength of approximately $\sigma(-2)\approx0.119$. This design avoids assigning artificial sequential meaning to mixed drug tokens and leaves BiIntention responsible for explicit cross-modal matching.

\subsection{Drug-Target Interaction Modeling}

Drug and target representations are fused by BiIntention, which serves as the explicit cross-modal matching module. We keep BiIntention in the main model because drug tokens mix SMILES semantics, graph topology, and conformer geometry, so concatenating drug and protein tokens into a single sequence is not a natural fit for state-space propagation. In this design, Mamba refines the ordered protein stream, while BiIntention remains responsible for explicit drug-target matching.

\subsection{Prediction Head and Training Objective}

The fused interaction representation is passed to an MLP classifier that outputs a logit for the interaction label. All experiments optimize binary cross-entropy with logits and use positive-class reweighting derived from the training split. Some historical runs included a counterfactual regularization branch, but in the current BiIntention implementation the attention output required by that branch is unavailable, so the counterfactual loss is effectively inactive. We therefore do not treat counterfactual regularization as a contribution in this draft.

\subsection{Algorithm Procedure}

For each drug-target pair, the model first loads cached multimodal drug features and protein sequence features. It then computes target-aware conformer aggregation, ligand-conditioned protein modulation, protein-side Mamba residual enhancement, BiIntention interaction fusion, and MLP prediction. During training, the prediction logit is optimized with weighted binary cross-entropy; during evaluation, the decision threshold is selected on the validation split and applied unchanged to the test split.

\subsection{Chapter Summary}

This chapter defines TAMR-DTI as a target-aware multimodal DTI model. The two paper-facing method components are target-aware conformer modulation and protein state-space enhancement before BiIntention. The next section evaluates the resulting full model on four benchmark DTI datasets.
